\chapter{무게가 다른 화살표}

지금까지 화살표는 모두 똑같았어요.
하지만 실제 기록은 그렇지 않아요.
스티커 1장을 보낸 것과 100장을 보낸 것은 다르지요.
어제 보낸 것과 일 년 전에 보낸 것도 달라요.

\section{굵은 화살표}

화살표마다 \textbf{무게}를 적어 둘 수 있어요.
하준이가 서연이에게 스티커 9장을, 도윤이에게 1장을 보냈다면
서연이 쪽 화살표에는 9, 도윤이 쪽 화살표에는 1을 적어요.

쪽지 놀이도 조금 바꿔요.
하준이가 “화살표를 따라가” 카드를 뽑으면, 열 번 가운데 아홉 번은 서연이에게, 한 번은 도윤이에게 쪽지를 넘겨요.
물 붓기라면 하준이의 물 가운데 90\%가 서연이에게, 10\%가 도윤이에게 가요.
무거운 화살표로 물이 더 많이 흐르는 거예요.

\section{오래된 일은 조금 가볍게}

일 년 전에 친구에게 스티커를 많이 준 일이 있다고 해 봐요.
그 일은 지금도 중요할까요?
조금은 중요하지만, 어제 일만큼은 아니겠지요.
사람의 기억도 시간이 지나면 흐려지고요.

그래서 화살표의 무게에 \textbf{신선도}를 곱해요.
신선도는 오늘 일어난 일이면 거의 1에 가깝고, 시간이 지날수록 줄어들어요.
제 논문에서 쓴 신선도는 이렇게 생겼어요.

\begin{figure}[h]
\centering
\begin{tikzpicture}
\begin{axis}[
  width=0.92\linewidth, height=5.2cm,
  xmin=0, xmax=400, ymin=0, ymax=1.05,
  xlabel={얼마나 오래된 일인가 (일)}, ylabel={신선도},
  xtick={0,90,180,270,365}, ytick={0,0.5,1},
  grid=major, grid style={gray!25},
  tick label style={font=\scriptsize}, label style={font=\small},
]
\addplot[domain=0:400, samples=120, line width=1.4pt, color=sendcol] {1/(1+exp(0.01*(x-180)))};
\addplot[only marks, mark=*, mark size=2pt, color=inkblue] coordinates {(0,0.858) (90,0.711) (180,0.5) (365,0.136)};
\node[font=\scriptsize, anchor=south west] at (axis cs:2,0.86) {오늘: 0.86};
\node[font=\scriptsize, anchor=south west] at (axis cs:92,0.72) {석 달 전: 0.71};
\node[font=\scriptsize, anchor=north east] at (axis cs:178,0.49) {반년 전: 0.5};
\node[font=\scriptsize, anchor=north east] at (axis cs:362,0.11) {일 년 전: 0.14};
\end{axis}
\end{tikzpicture}
\caption{시간이 지날수록 줄어드는 신선도. 반년(180일) 전 일은 무게가 절반이 돼요.}
\end{figure}

반년 전 일은 무게가 절반, 일 년 전 일은 7분의 1쯤 돼요.
그래서 오래전에 큰 송금을 했더라도 요즘 아무 일도 없으면 점수가 조금씩 내려가요.

\section{챔피언: 송금 랭크}

이제 지금까지 만든 것을 모두 합쳐 볼게요.

\begin{enumerate}
\item 송금 줄로 지도를 그린다.
\item 화살표의 무게는 보낸 양에 신선도를 곱한 값이다.
\item 그 지도에서 쪽지 돌리기 놀이(페이지랭크)를 한다.
\end{enumerate}

이 방법에는 이름이 있어요.
2023년에 열린 한 국제 학회에서 도, 도, 응우옌이라는 세 연구자가 블록체인 지갑의 평판을 재려고 발표한 방법이에요.
우리말로 옮기면 “적응형 가중 페이지랭크”이고, 논문에서는 영어 이름의 머리글자를 따서 AWP라고 불러요.
이 책에서는 송금 지도에서 하는 페이지랭크라는 뜻으로 \textbf{송금 랭크}라고 부를게요.
송금 기록으로 지갑의 평판을 재는, 이미 나와 있는 방법이지요.

제 논문에서 송금 랭크는 \emph{챔피언}이에요.
새로 만든 방법이 쓸모 있다고 말하려면 챔피언과 겨루어야 하니까요.
권투 경기에 비유하면, 저는 도전자를 데리고 챔피언에게 도전하러 가는 셈이에요.
도전자가 누구인지는 다음 장에서 소개할게요.

\section{챔피언을 정직하게 모셔 오기}

여기서 고백할 것이 하나 있어요.
제 논문의 송금 랭크는 원래 논문에 나온 모습과 두 군데가 달라요.

\begin{itemize}
\item 원래 송금 랭크는 보낸 양을 그대로 쓰지 않고, 아주 큰 송금도 “한 번”쯤으로 줄여서 써요.
      제 논문의 송금 랭크는 보낸 양을 그대로 썼어요.
\item 원래 송금 랭크에서는 선생님이 쪽지를 아무에게나 주지 않고, 요즘 송금을 활발히 한 친구에게 더 자주 줘요.
      제 논문의 송금 랭크에서는 선생님이 아무에게나 똑같이 줬어요.
\end{itemize}

도전자와 겨룰 챔피언을 마음대로 바꾸면 공정하지 않겠지요.
그래서 저는 원래 논문에 나온 그대로의 송금 랭크도 따로 만들어서 같은 시험을 한 번 더 치렀어요.
결과가 달라지는 곳은 13장과 14장에서 그대로 알려 줄게요.

\begin{deeper}[신선도와 원래 송금 랭크의 식]
신선도는 로지스틱 함수 $\sigma(\Delta t) = 1/\bigl(1+e^{k(\Delta t - t_0)}\bigr)$이고,
논문에서는 $k=0.01$(하루당), $t_0=180$일을 썼어요.
$\Delta t$는 점수를 매기는 날을 기준으로 그 송금이 며칠 전에 일어났는지예요.

원래 송금 랭크(도, 도, 응우옌, 2023년)에서 송금 하나의 무게는 $T(t)\times V(z)$이고,
$V(z) = 2/(1+e^{-bz}) - 1$은 보낸 양 $z$를 0과 1 사이로 줄이는 함수예요.
원래 논문은 $b$ 값을 정해 두지 않았어요.
토큰의 양은 아주 큰 정수로 기록되기 때문에 $b$가 1 근처라면 아주 작은 송금이 아닌 한 $V(z)$는 거의 1이 돼요.
그러면 송금 한 번이 오래된 정도만큼 줄어든 채 “한 번”으로 세어져요.
되돌아가는 몫(선생님 몫)은 각 지갑이 최근에 얼마나 활발히 보냈는지에 비례해서 나누어요.
\end{deeper}

\begin{learned}
\begin{itemize}
\item 화살표에 무게를 주면, 무거운 화살표로 쪽지가 더 자주 가요.
\item 오래된 일은 신선도를 곱해서 가볍게 쳐요. 반년 전 일은 절반이에요.
\item 송금 지도에 이 두 가지를 더해 페이지랭크를 한 것이 송금 랭크예요. 송금 랭크가 이 논문의 챔피언이에요.
\item 공정하게 겨루려고, 원래 논문 그대로의 송금 랭크와도 따로 겨루었어요.
\end{itemize}
\end{learned}
